\clearpage
\VerdierPhysicalPageAnchor{0163}%
\begingroup\small\setlength{\parskip}{0pt}\setlength{\emergencystretch}{1.5em}
\VerdierSemanticAnchor{ch4-4-1-2-b}
b) Un foncteur de $\mathcal{F}\ell_2(\mathcal J)$ dans la catégorie des
triangles distingués de $\mathcal D$, de la forme
\[
\omega_2=(\mathord\bullet\xrightarrow f\mathord\bullet\xrightarrow g\mathord\bullet)
\longmapsto
X_f\xrightarrow{X*u_1(\omega_2)}X_{gf}
\xrightarrow{X*u_2(\omega_2)}X_g
\xrightarrow{\delta(\omega_2)}X_f[1].
\]
Pour désigner les objets spectraux de type $\mathcal J$, nous utiliserons la
notation
\[
(X_f, f\in\operatorname{Ob}\mathcal{F}\ell_1(\mathcal J),\
\delta(\omega_2),
\ \omega_2\in\operatorname{Ob}\mathcal{F}\ell_2(\mathcal J)),
\]
ou plus simplement $(X_f,\delta)$ quand aucune confusion n'en résulte.

\VerdierSemanticAnchor{ch4-4-1-3}
\medskip\noindent\textbf{4.1.3.}\quad
Soient $C$ une catégorie additive et
$F\colon\mathcal J\longrightarrow\operatorname{comp}(C)$ un foncteur. On a
montré au chapitre I, (3.3.6), comment associer à $F$ un objet spectral de
type $\mathcal J$ à valeurs dans $K(C)$. L'objet spectral décrit au chapitre
I, (3.3.6), est appelé l'\emph{objet spectral associé au foncteur $F$}.

\VerdierSemanticAnchor{ch4-4-1-4}
\medskip\noindent\textbf{4.1.4.}\quad
Soient $\mathcal A$ une catégorie abélienne et $\mathcal J$ une catégorie.
Un \emph{objet spectral de type $\mathcal J$ à valeurs dans $\mathcal A$}
est constitué par :

a) Une suite de foncteurs
$H^n\colon\mathcal{F}\ell_1(\mathcal J)\longrightarrow\mathcal A$,
$n\in\mathbb Z$.

b) La donnée, pour tout objet
$\omega_2=(\mathord\bullet\xrightarrow f\mathord\bullet\xrightarrow g\mathord\bullet)$
de $\mathcal{F}\ell_2(\mathcal J)$ et tout entier $n$, d'un morphisme
$\delta^n(\omega_2)\colon H^n(g)\to H^{n+1}(f)$, fonctoriel en $\omega_2$,
tel que le diagramme
\[
\cdots\longrightarrow H^n(f)
\xrightarrow{H^n*u_1(\omega_2)}H^n(gf)
\xrightarrow{H^n*u_2(\omega_2)}H^n(g)
\xrightarrow{\delta^n(\omega_2)}H^{n+1}(f)
\longrightarrow\cdots
\]
soit une suite exacte illimitée de $\mathcal A$.

Nous utiliserons la notation
\[
(H^n(f), f\in\operatorname{Ob}\mathcal{F}\ell_1(\mathcal J),\
\delta^n(\omega_2),
\ \omega_2\in\operatorname{Ob}\mathcal{F}\ell_2(\mathcal J))
\]
pour désigner les objets spectraux de type $\mathcal J$, ou plus simplement
$(H^n(f),\delta)$ quand aucune confusion n'en résulte.

\VerdierSemanticAnchor{ch4-4-1-5}
\medskip\noindent\textbf{4.1.5.}\quad
Soit $F\colon\mathcal D\to\mathcal D'$ un foncteur exact entre deux
catégories triangulées. Soit $(X_f,\delta)$ un objet spectral de type
$\mathcal J$ à valeurs dans $\mathcal D$. Le couple de foncteurs
$(FX_f,F\delta)$ est un objet spectral de type $\mathcal J$ à valeurs dans
$\mathcal D'$. De même, un foncteur cohomologique
$H\colon\mathcal D\to\mathcal A$ d'une catégorie triangulée dans une
catégorie abélienne transforme les objets spectraux à valeurs dans
$\mathcal D$ en objets spectraux à valeurs dans $\mathcal A$.

\VerdierSemanticAnchor{ch4-4-1-6}
\medskip\noindent\textbf{4.1.6.}\quad
Soient $\mathcal D$ une catégorie triangulée et $(X_f,\delta)$,
$(Y_f,\delta')$ deux objets spectraux de type $\mathcal J$ à valeurs dans
$\mathcal D$. Un \emph{morphisme d'objets spectraux} est un morphisme de
foncteurs
\[
m_f\colon X_f\longrightarrow Y_f,\qquad
f\in\operatorname{Ob}\mathcal{F}\ell_1(\mathcal J),
\]
\endgroup
